%% SCRAP: architecture/03-architecture/physics-engine/ssm-equations
%% SOURCE: docs/working/architecture/03-architecture/physics-engine/ssm-equations.md
%% STATUS: WORKING
%% FITS: dev-guide/ch-physics
%% EDITORIAL: lifted — prose rewritten to press voice

\section{Mathematical Models of the Steady-State Machine}

%% PATENT: This scrap covers patent-adjacent empirical relationships (the capacity
%% law and its candidate claim language). Claim wording in the source is summarized
%% as description only; no claim language is drafted here. Flag for Bob before any
%% promotion that touches the claims.

This section presents performance models derived empirically from the
Steady-State Machine (SSM) experimental record: 51{,}840 runs of the StarForth
VM across 128 feedback-loop configurations. Five relationships were extracted
from the data with measurable precision. The headline result is a capacity law,
$\Lambda(\mathrm{DoF}) = 4096 / (\mathrm{DoF}+1)$, which holds with a 0.00\%
coefficient of variation across all stable configurations.

These models are exploratory. Several borrow functional forms and names from
physics --- the Lorentz factor, the Friedmann equation, the Schwarzschild
radius --- purely because the data happens to fit those forms. The physics
vocabulary is a modeling convenience for describing adaptive-system behavior; it
is not a claim that the runtime obeys relativity or cosmology. Throughout,
\textit{degrees of freedom} (DoF) denotes the number of enabled feedback loops,
ranging from 0 to 7.

\subsection{Equation 1 --- Performance Scaling with Load}

Execution time scales with the number of enabled loops according to a Lorentz-like
factor:

\begin{equation}
\tau(\mathrm{DoF}) = \tau_0 \cdot \gamma(\mathrm{DoF}),
\qquad
\gamma = \frac{1}{\sqrt{1 - \beta^2}},
\qquad
\beta^2(\mathrm{DoF}) = 0.0736\,\mathrm{DoF} + 0.1464,
\end{equation}

with $\tau_0 = 31.67$\,ms/word the rest-state performance. The empirical fit is
$R^2 = 0.938$. The $\beta^2$ term grows linearly with DoF, so each additional
loop carries increasing marginal overhead. The structural similarity to
relativistic time dilation is offered as a predictive model, not a physical
claim.

\subsection{Equation 2 --- Window Capacity Law $\Lambda$}

The central empirical finding is an inverse capacity law:

\begin{equation}
\Lambda(\mathrm{DoF}) = \frac{W_0}{\mathrm{DoF} + 1},
\qquad W_0 = 4096.
\end{equation}

Equivalently, the product $\Lambda \times (\mathrm{DoF}+1) = W_0$ is conserved
across all stable configurations.

\begin{table}[ht]
\centering
\small
\begin{tabular}{rrrr}
\toprule
DoF & Ideal $\Lambda$ & $\Lambda \times (\mathrm{DoF}+1)$ & Stability (CV) \\
\midrule
0 & 4096.0 & 4096.0 & 17.34\% \\
1 & 2048.0 & 4096.0 & 16.70\% \\
2 & 1365.3 & 4096.0 & 15.31\% \\
3 & 1024.0 & 4096.0 & 13.91\% \\
4 & 819.2  & 4096.0 & 13.77\% \\
5 & 682.7  & 4096.0 & 13.96\% \\
6 & 585.1  & 4096.0 & 14.33\% \\
7 & 512.0  & 4096.0 & 22.06\% \\
\bottomrule
\end{tabular}
\caption{Window capacity verification. The conserved product holds at $4096.0 \pm 0.0$, a 0.00\% coefficient of variation.}
\end{table}

The symbol $\Lambda$ is borrowed from cosmology for convenience: it denotes the
effective window capacity per degree of freedom, with $W_0 = 4096$ bytes an
empirically determined constant. As loops are added, proportional window
capacity must fall to preserve stability, and the conserved product behaves as a
stability invariant. Notably, $W_0 = 4096 = 2^{12}$ was not designed --- it
emerged from experimental optimization, and may reflect memory page-size
alignment, cache-line constraints, or buffer-size effects.

\subsection{Equation 3 --- Expansion Dynamics (Friedmann Form)}

The growth in execution time as loops activate fits a Friedmann-like form:

\begin{equation}
H^2(\mathrm{DoF}) \approx c_1\,\rho(\mathrm{DoF}) + c_2\,\Lambda(\mathrm{DoF}) + c_3,
\end{equation}

with the fitted relation

\begin{equation}
H^2 = -2.24\times10^{12}\,\mathrm{DoF} \;-\; 4.51\times10^{9}\,\Lambda \;+\; 2.11\times10^{13}.
\end{equation}

Here $H$ stands in for the performance growth rate, $\rho$ for the degrees of
freedom, and $\Lambda$ for window-capacity pressure. The negative coefficients
indicate that both DoF and $\Lambda$ contribute to deceleration --- the system
resists unbounded growth through feedback regulation.

\subsection{Equation 4 --- Instability Threshold (Schwarzschild Form)}

A Schwarzschild-like radius models the onset of instability:

\begin{equation}
r_s(\mathrm{DoF}) = \frac{\mathrm{DoF}}{W_{\mathrm{norm}}},
\qquad W_{\mathrm{norm}} = \frac{\text{window}}{4096}.
\end{equation}

In this analogy $r_s$ is the instability threshold, DoF the computational load,
and $W_{\mathrm{norm}}$ the available space. The configurations with the highest
coefficient of variation sit near this threshold.

\begin{table}[ht]
\centering
\small
\begin{tabular}{rrrrr}
\toprule
Config & DoF & $r_s$ & CV & Performance \\
\midrule
1  & 1 & 1.00 & 26.90\% & 32.43\,ms \\
84 & 3 & 3.00 & 26.36\% & 34.94\,ms \\
46 & 4 & 4.00 & 25.70\% & 50.27\,ms \\
59 & 5 & 5.00 & 25.31\% & 50.43\,ms \\
\bottomrule
\end{tabular}
\caption{Configurations nearest the instability threshold. Where $r_s \ge \mathrm{DoF}$, the system tends toward instability.}
\end{table}

\subsection{Equation 5 --- Energy--Momentum Form}

Total execution time fits an energy--momentum relation:

\begin{equation}
E^2 = E_0^2 + p^2,
\qquad E = \tau(\mathrm{DoF}),
\qquad E_0 = \tau_0 = 31.67\,\mathrm{ms},
\qquad p(\mathrm{DoF}) = 5.20\,\mathrm{DoF} + 6.55.
\end{equation}

Here $E_0$ is the all-loops-off baseline, $p$ the additional time cost from
loops, and $E$ total execution time. Momentum scales linearly with DoF
($R = 0.425$, $p < 10^{-6}$).

\begin{table}[ht]
\centering
\small
\begin{tabular}{rrrrr}
\toprule
Config & DoF & $E_{\text{total}}$ & $E_{\text{rest}}$ & $p$ \\
\midrule
0   & 0 & 31.67 & 31.67 & 0.00 \\
35  & 3 & 31.59 & 31.67 & 0.00 \\
55  & 5 & 31.91 & 31.67 & 3.88 \\
124 & 5 & 59.48 & 31.67 & 50.35 \\
\bottomrule
\end{tabular}
\caption{Energy--momentum decomposition for selected configurations.}
\end{table}

Two configurations are notable. Config~\#35 carries three enabled loops yet sits
at near-zero momentum, suggesting an optimal arrangement where loop interactions
cancel. Config~\#124, the worst performer, carries the same DoF as Config~\#55
but enormous momentum (50.35) --- a destructive combination of loops.

\subsection{Interpretation and Open Experiments}

The recurring observation is that these adaptive-runtime relationships fit
functional forms also found in general relativity, thermodynamics, and quantum
mechanics. The strongest of them, the capacity law, is exact to within
measurement. Whether this reflects a deep property of feedback-driven systems or
an emergent coincidence remains open; the SSM data is presented as evidence that
the functional forms are useful predictive models.

Several follow-on experiments are proposed:

\begin{itemize}
  \item \textbf{Window-size scaling.} Vary the window across 1024, 2048, 4096,
        8192, and 16384 bytes and test whether $\Lambda(\mathrm{DoF}) =
        W/(\mathrm{DoF}+1)$ holds, i.e.\ whether the 4096 constant scales
        linearly with window size.
  \item \textbf{Critical window capacity.} Reduce the window in steps and locate
        a $W_{\text{crit}}$ below which the selector collapses to the all-off
        configuration.
  \item \textbf{Maximum degrees of freedom.} Extend beyond seven loops and test
        whether the capacity law holds up to a $\mathrm{DoF}_{\max}$, with the
        candidate $\mathrm{DoF}_{\max} \approx \log_2(W_0) = 12$.
\end{itemize}

%% PATENT: The source proposes patent claims based on the capacity law, the
%% Lorentz-like scaling, and the Schwarzschild-like threshold, and outlines
%% publication and funding strategy. That material is intentionally omitted here
%% pending Bob's instruction; do not draft claim language in this scrap.
%% TODO(bob): confirm which empirical relationships may appear in citable form.
